% FILE: 00_project_identity.tex
% PROJECT: experiments-core
% THESIS ROLE: benchmark-experiments-and-equation-verification

\section{Project Identity}
\label{sec:experiments-core-identity}

\subsection{Purpose}
\label{sec:experiments-core-identity-purpose}

The \texttt{experiments/} directory serves as the empirical verification layer of the
process-intelligence research program. Its purpose is threefold: to mathematically verify
the foundational equations that govern conformance checking; to provide executable benchmark
fixtures that compare wasm4pm against PM4Py across performance, type-safety, and
zero-knowledge dimensions; and to demonstrate the full evidence chain from raw OCEL logs
through typed Admission gates to board-admissible conformance claims. The experiments do
not simulate conformance---they manufacture it from lawful primitives and record cryptographic
receipts that a board-level claim may cite without further interpretation. Every claim made
in downstream manufacturing stages must trace to one of the experiment artifacts housed
here or be declared AUTHOR THESIS.

\subsection{Repository Surface}
\label{sec:experiments-core-identity-surface}

The repository surface comprises three distinct layers. The first layer consists of
executable Python validation scripts: \texttt{validate\_log\_fitness.py},
\texttt{validate\_ocpq\_refinement.py}, and \texttt{verify\_equations.py}. These scripts
run without mocks and exit with assertion failures on any incorrect result. The second
layer is a structured set of Markdown experiment samples with concrete JSON schemas covering
18 domains: Petri net token replay, OCEL 2.0 lifecycle, cryptographic replay receipts,
OTel-BLAKE3 integration, raw laundering refusal, and M\&A claim mappings, among others.
The third layer is a fully interactive browser-based process-mining dashboard
(\texttt{visualizer/}) that implements A* alignment, EWMA drift detection, a SHA-256
audit ledger, LTL Declare constraint checking, and Petri net token game animation. Two
subdirectories---\texttt{ma-deck-projection/} and \texttt{wasm4pm-gap-analysis/}---are
currently empty and represent deferred manufacturing work (see Section~\ref{sec:experiments-core-open}).

\subsection{Primary Research Role}
\label{sec:experiments-core-identity-role}

Within the dissertation, this project occupies the role of \emph{benchmark experiments
and equation verification}. It provides the only executable proofs that the Van der Aalst
weighted log fitness equation and the OCPQ query variable refinement relation hold on
concrete data. It further provides the only reproducible performance comparison between
PM4Py and wasm4pm. Without this layer, claims about conformance speed, type safety, and
zero-knowledge compatibility remain conjectural. The experiments directory is therefore
a prerequisite for any board-admissible claim that cites fitness metrics, alignment costs,
or ZKP cycle counts.

\subsection{Key Artifacts}
\label{sec:experiments-core-identity-artifacts}

The key artifacts are: (1)~\texttt{verify\_equations.py}---a single gate that runs
five fitness cases and four OCPQ cases and reports 100\% correctness; (2)~\texttt{audit\_\_experiment\_completeness.md}---records 18/18 experiments as PASSED with concrete JSON schemas and zero stubs;
(3)~\texttt{checkpoint\_\_experiments\_complete.md}---declares VERIFIED \& COMPLETE;
(4)~\texttt{EVIDENCE\_CHAIN\_TRACE.md}---documents the 7-step evidence chain from raw OCEL
through receipt to board claim; (5)~\texttt{visualizer/bindings.d.ts}---TypeScript type
projections of the Rust typestate system crossing the WASM boundary; and (6)~the
\texttt{pm4py-comparison/} benchmark reports recording DFG Discovery 37.5$\times$,
Alpha Miner 34.2$\times$, and Token Replay 39.3$\times$ speedups for wasm4pm over PM4Py.
